\documentclass[12pt]{article}
\usepackage[utf8]{inputenc}
\usepackage[T1]{fontenc}
\usepackage{hyperref}
\title{Fataleka Language Resources for K-12 Curriculum Development}
\author{Compiled by Claude Sonnet 5.5 (claude-sonnet-5-5) for Andr\'as Kornai}
\date{}
\begin{document}
\maketitle
\section{Introduction}
Fataleka (ISO 639-3 far) is a Southeast Solomonic Austronesian language of Malaita, Solomon Islands, with roughly 6,700 speakers. Very little educational material exists; this report lists only what was found in web searches.
\section{Primary Educational Resources}
No primers, textbooks or school materials were found in search results.
\section{Linguistic and Reference Materials}
No dictionary or grammar was found. A short verb list is on Verbix: \url{https://docs.verbix.com/Languages/Fataleka}. General overview: \url{https://en.wikipedia.org/wiki/Fataleka_language}. Search results also pointed to Arthur Capell materials held by Paradisec, but no URL was seen, so none is cited.
\section{Religious and Community Resources}
A Fataleka New Testament (\emph{Farongoe Lea Ana Fataleka}, 2017, Wycliffe/SITAG) is available at \url{https://dbs.org/bibles/FARWBT}. People group profile: \url{https://joshuaproject.net/people_groups/11192}.
\section{Digital and Audio Resources}
Global Recordings Network lists Fataleka audio programmes (Bible lessons, stories, songs): \url{https://globalrecordings.net/en/language/3617}. The New Testament also has an audio version through the DBS page above.
\section{Curriculum Development Context}
The language is not known to be taught in schools. A native-language curriculum would have to be built almost from scratch, using the New Testament as the main written corpus, and would need new orthographic and lexical documentation.
\end{document}